\section{Discussion}

CiTT reframes AI tutoring for engineering courses as a model-grounding problem. A tutor can ask helpful questions and explain concepts, but for circuits and biomedical instrumentation, the strongest teaching moment often comes from seeing where a quantity is measured, how a model behaves, and which assumptions are missing. CiTT therefore treats LLM output as one component in a broader MATLAB-centered workflow.

The evidence also shows that LLM-only tutoring should not be dismissed. The Gemini-only baseline performed well on the RC arithmetic and gave plausible TEVC intuition. The limitation is not language itself; it is the absence of executable, inspectable artifacts. CiTT's advantage grows as tasks require topology, symbolic caveats, probes, transients, ADC behavior, saturation, or fault reasoning.

For educational deployment, CiTT is most natural in courses and labs already using MATLAB/Simulink/Simscape. The distribution format is a MATLAB toolbox package or source zip. Costs are mostly software access, setup time, API usage, and maintenance rather than hardware manufacturing.

\begin{table*}[t]
\centering
\small
\begin{tabular}{p{0.22\textwidth}p{0.40\textwidth}p{0.28\textwidth}}
\toprule
BMES criterion & Evidence in CiTT & Repository artifact \\
\midrule
Product need and market potential & BME courses need circuit/instrumentation tutoring that connects diagrams to executable models. & \texttt{docs/release\_setup.md}, live evidence report \\
Utility and novelty & MATLAB-native agentic tutor with Simscape artifacts, focus maps, probes, and evidence export. & \texttt{matlab/}, \texttt{submission\_assets/live\_gui\_evidence/} \\
Technical feasibility & Three live benchmarks and installed \texttt{.mltbx} smoke/examples verification. & \texttt{release/package\_log.md}, \texttt{release/example\_repro/} \\
Budget/economic plan & Software package leveraging academic MATLAB infrastructure and API usage. & BMES application draft \\
Limitations and risk & Educational-only scope, symbolic caveats, non-settling/saturation evidence, pending Lab Delta CSV. & Evidence report, limitations raw material \\
\bottomrule
\end{tabular}
\caption{Mapping between BMES judging criteria and repository evidence.}
\label{tab:bmes_mapping}
\end{table*}

